% GENERATED FILE. Edit canonical lessons, then run npm run generate:books.
% canonical-source-hash: fnv1a64:df459ac7c8957a96
% canonical-lessons: MR-C437-garva, MR-C437-himvarsav, MR-C437-jhombne, MR-C437-gothne, MR-C437-vitalne

\chapter{A Biting Chill}
\label{ch:mr-a-biting-chill}
\hlchaptermodality{437}

\begin{chapteropening}
\textbf{By the end of this chapter:} \emph{I can talk about coolness and snowfall, and say the wind nipped, water froze and snow melted.}

Complete the last lesson of chapter 437: I can talk about coolness and snowfall, and say the wind nipped, water froze and snow melted.
\end{chapteropening}

\section[\texorpdfstring{\mr{गारवा} (gārvā)}{गारवा (gārvā)}]{\mr{गारवा} (gārvā) — coolness; a chill}

\label{lesson:MR-C437-garva}



\begin{quote}
Before the new one: say the Marathi for humid, then the Marathi for to blaze.
\end{quote}

\subsection*{You'll want to know: \mr{गारवा}}
\textbf{\mr{गारवा}} — \emph{gārvā} — \textquotedblleft{}coolness; a chill\textquotedblright{}.

\begin{quote}
\mr{संध्याकाळी} \mr{हवेत} \mr{गारवा} \mr{आला}.
\end{quote}

\begin{quote}
\emph{sandhyākāḷī havet gārvā ālā.}
\end{quote}

\begin{quote}
In the evening a chill came into the air.
\end{quote}

Masculine.

A pleasant coolness; \mr{गारवा} \mr{येणे}.

\begin{cousinweb}
\textbf{\mr{गारवा}} (\emph{gārvā}) is built on \emph{gār}, \textquotedblleft{}cool, cold\textquotedblright{}, with a noun ending; its origin is not settled.
\end{cousinweb}

\begin{grammarlens}[title={What you've built}]
One more word for getting around.
\end{grammarlens}

\subsection*{Your turn}
\begin{itemize}
\raggedright
  \item \emph{Say it:} \emph{gārvā}
  \item \emph{Say it:} \emph{gārvā}, once more
  \item \emph{Say it:} \emph{sandhyākāḷī havet gārvā ālā.} — In the evening a chill came into the air.
  \item \emph{Recall:} say the Marathi for humid, then the Marathi for to blaze, then say \emph{gārvā} again
\end{itemize}

\begin{tcolorbox}[breakable,colback=teal!4,colframe=teal!35!black,title={Before you move on}]
What does \mr{गारवा} mean? (\textquotedblleft{}Coolness; a chill\textquotedblright{}.) Say it once more.
\end{tcolorbox}
\section[\texorpdfstring{\mr{हिमवर्षाव} (himvarṣāv)}{हिमवर्षाव (himvarṣāv)}]{\mr{हिमवर्षाव} (himvarṣāv) — a snowfall}

\label{lesson:MR-C437-himvarsav}



\begin{quote}
Before the new one: say the Marathi for to blaze, then the Marathi for coolness.
\end{quote}

\subsection*{You'll want to know: \mr{हिमवर्षाव}}
\textbf{\mr{हिमवर्षाव}} — \emph{himvarṣāv} — \textquotedblleft{}a snowfall\textquotedblright{}.

\begin{quote}
\mr{डोंगरावर} \mr{रात्रभर} \mr{हिमवर्षाव} \mr{झाला}.
\end{quote}

\begin{quote}
\emph{ḍoṅgarāvar rātrabhar himvarṣāv jhālā.}
\end{quote}

\begin{quote}
It snowed on the mountain all night.
\end{quote}

Masculine.

\mr{हिम}, snow, and \mr{वर्षाव}, a shower.

\begin{cousinweb}
\textbf{\mr{हिमवर्षाव}} (\emph{himvarṣāv}) is built from Sanskrit \emph{hima}, \textquotedblleft{}snow\textquotedblright{} (Monier-Williams), and \textbf{\mr{वर्षाव}}, \textquotedblleft{}a shower\textquotedblright{}, from the family of \emph{varṣati}, \textquotedblleft{}rains\textquotedblright{} (CDIAL 11394).
\end{cousinweb}

\begin{grammarlens}[title={What you've built}]
One more word for getting around.
\end{grammarlens}

\subsection*{Your turn}
\begin{itemize}
\raggedright
  \item \emph{Say it:} \emph{himvarṣāv}
  \item \emph{Say it:} \emph{himvarṣāv}, once more
  \item \emph{Say it:} \emph{ḍoṅgarāvar rātrabhar himvarṣāv jhālā.} — It snowed on the mountain all night.
  \item \emph{Recall:} say the Marathi for to blaze, then the Marathi for coolness, then say \emph{himvarṣāv} again
\end{itemize}

\begin{tcolorbox}[breakable,colback=teal!4,colframe=teal!35!black,title={Before you move on}]
What does \mr{हिमवर्षाव} mean? (\textquotedblleft{}A snowfall\textquotedblright{}.) Say it once more.
\end{tcolorbox}
\section[\texorpdfstring{\mr{झोंबणे} (jhombṇe)}{झोंबणे (jhombṇe)}]{\mr{झोंबणे} (jhombṇe) — to sting; to nip}

\label{lesson:MR-C437-jhombne}



\begin{quote}
Before the new one: say the Marathi for coolness, then the Marathi for a snowfall.
\end{quote}

\subsection*{You'll want to know: \mr{झोंबणे}}
\textbf{\mr{झोंबणे}} — \emph{jhombṇe} — \textquotedblleft{}to sting; to nip\textquotedblright{}.

\begin{quote}
\mr{पहाटेचा} \mr{थंड} \mr{वारा} \mr{अंगाला} \mr{झोंबला}.
\end{quote}

\begin{quote}
\emph{pahāṭecā thaṇḍ vārā aṅgālā jhomblā.}
\end{quote}

\begin{quote}
The cold dawn wind stung my body.
\end{quote}

Past (he …): \mr{तो} \mr{झोंबला}, \emph{to jhomblā}.

What is stung takes -\mr{ला}: \mr{अंगाला}, \mr{कानाला} \mr{झोंबणे}.

\begin{cousinweb}
\textbf{\mr{झोंबणे}} (\emph{jhombṇe}) is a Marathi verb for cold, smoke or spice that stings; its origin is not settled.
\end{cousinweb}

\begin{grammarlens}[title={What you've built}]
One more word for getting around.
\end{grammarlens}

\subsection*{Your turn}
\begin{itemize}
\raggedright
  \item \emph{Say it:} \emph{jhombṇe}
  \item \emph{Say it:} \emph{jhombṇe}, once more
  \item \emph{Say it:} \emph{pahāṭecā thaṇḍ vārā aṅgālā jhomblā.} — The cold dawn wind stung my body.
  \item \emph{Recall:} say the Marathi for coolness, then the Marathi for a snowfall, then say \emph{jhombṇe} again
\end{itemize}

\begin{tcolorbox}[breakable,colback=teal!4,colframe=teal!35!black,title={Before you move on}]
What does \mr{झोंबणे} mean? (\textquotedblleft{}To sting; to nip\textquotedblright{}.) Say it once more.
\end{tcolorbox}
\section[\texorpdfstring{\mr{गोठणे} (goṭhṇe)}{गोठणे (goṭhṇe)}]{\mr{गोठणे} (goṭhṇe) — to freeze}

\label{lesson:MR-C437-gothne}



\begin{quote}
Before the new one: say the Marathi for a snowfall, then the Marathi for to sting.
\end{quote}

\subsection*{You'll want to know: \mr{गोठणे}}
\textbf{\mr{गोठणे}} — \emph{goṭhṇe} — \textquotedblleft{}to freeze\textquotedblright{}.

\begin{quote}
\mr{थंडीत} \mr{तळ्याचं} \mr{पाणी} \mr{गोठलं}.
\end{quote}

\begin{quote}
\emph{thaṇḍīt taḷyācaṁ pāṇī goṭhlaṁ.}
\end{quote}

\begin{quote}
In the cold the water of the pond froze.
\end{quote}

Past (he …): \mr{तो} \mr{गोठला}, \emph{to goṭhlā}.

Water and blood \mr{गोठतात}; not the same word as \mr{गोठा}, a cowshed.

\begin{cousinweb}
\textbf{\mr{गोठणे}} (\emph{goṭhṇe}) is a Marathi verb for water turning to ice; its origin is not settled.
\end{cousinweb}

\begin{grammarlens}[title={What you've built}]
One more word for getting around.
\end{grammarlens}

\subsection*{Your turn}
\begin{itemize}
\raggedright
  \item \emph{Say it:} \emph{goṭhṇe}
  \item \emph{Say it:} \emph{goṭhṇe}, once more
  \item \emph{Say it:} \emph{thaṇḍīt taḷyācaṁ pāṇī goṭhlaṁ.} — In the cold the water of the pond froze.
  \item \emph{Recall:} say the Marathi for a snowfall, then the Marathi for to sting, then say \emph{goṭhṇe} again
\end{itemize}

\begin{tcolorbox}[breakable,colback=teal!4,colframe=teal!35!black,title={Before you move on}]
What does \mr{गोठणे} mean? (\textquotedblleft{}To freeze\textquotedblright{}.) Say it once more.
\end{tcolorbox}
\section[\texorpdfstring{\mr{वितळणे} (vitaḷṇe)}{वितळणे (vitaḷṇe)}]{\mr{वितळणे} (vitaḷṇe) — to melt}

\label{lesson:MR-C437-vitalne}



\begin{quote}
Before the new one: say the Marathi for to sting, then the Marathi for to freeze.
\end{quote}

\subsection*{You'll want to know: \mr{वितळणे}}
\textbf{\mr{वितळणे}} — \emph{vitaḷṇe} — \textquotedblleft{}to melt\textquotedblright{}.

\begin{quote}
\mr{उन्हात} \mr{बर्फ} \mr{लवकर} \mr{वितळतो}.
\end{quote}

\begin{quote}
\emph{unhāt barpha lavkar vitaḷto.}
\end{quote}

\begin{quote}
Ice melts quickly in the sun.
\end{quote}

Past (he …): \mr{तो} \mr{वितळला}, \emph{to vitaḷlā}.

Ice, butter and wax \mr{वितळतात}.

\begin{cousinweb}
\textbf{\mr{वितळणे}} (\emph{vitaḷṇe}) is a Marathi verb for something solid turning liquid; its origin is not settled.
\end{cousinweb}

\begin{grammarlens}[title={What you've built}]
One more word for getting around.
\end{grammarlens}

\subsection*{Your turn}
\begin{itemize}
\raggedright
  \item \emph{Say it:} \emph{vitaḷṇe}
  \item \emph{Say it:} \emph{vitaḷṇe}, once more
  \item \emph{Say it:} \emph{unhāt barpha lavkar vitaḷto.} — Ice melts quickly in the sun.
  \item \emph{Recall:} say the Marathi for to sting, then the Marathi for to freeze, then say \emph{vitaḷṇe} again
\end{itemize}

\begin{tcolorbox}[breakable,colback=teal!4,colframe=teal!35!black,title={Before you move on}]
What does \mr{वितळणे} mean? (\textquotedblleft{}To melt\textquotedblright{}.) Say it once more.
\end{tcolorbox}
